% !PRINTED 606
\weberpage{606}
%<<CONT
% [page alignment corrected 2026-08-07] Printed 606 begins with this
% continuation from i0630 and ends after the §179 footnote.  The former
% "Wir setzen ..." tail is printed on 607 and already belongs to i0632.
so ist $f(x)$ irreducibel, und man hat in der Wahl der ganzen Zahlen $c$ noch Freiheit genug, um die rationalen Br\"uche $a_1, a_2 \ldots a_n$ einem beliebig gegebenen Werthsystem beliebig nahe zu bringen.

Damit ist aber der Satz bewiesen:

4. \gesperrt{Es giebt von jedem beliebigen Primzahlgrade unendlich viele Gleichungen mit rationalen Co\"efficienten ohne Affect.}

Der Beweis, der hier gef\"uhrt ist, zeigt, dass solche affectfreie Gleichungen gefunden werden k\"onnen, deren Co\"efficienten ein endliches Gebiet \"uberall dicht erf\"ullen, das Gebiet n\"amlich, in dem die reellen Co\"efficienten der Gleichungen mit nur zwei
conjugirt imagin\"aren Wurzeln liegen. Der Beweis ist also, abgesehen davon, dass er sich nur auf Primzahlgrade beschr\"ankt, auch insofern nicht ersch\"opfend, als er uns keinen Aufschluss giebt \"uber die \"ubrigen Gebiete, in denen aller Wahrscheinlichkeit nach die Sache sich ebenso verh\"alt.

\begin{center}
\S.\,179.

\gesperrt{Aufl\"osung durch reelle Radicale.}
\end{center}

Bei der Aufl\"osung der cubischen Gleichungen mit reellen Co\"efficienten hat man von Alters her die zwei F\"alle unterschieden, in denen die Discriminante negativ oder positiv ist. Im zweiten Falle hat man, obwohl die Gleichung dann gerade drei reelle Wurzeln hat, bei der Anwendung der \gesperrt{Cardanischen Formel} die dritte Wurzel aus einem imagin\"aren Ausdruck zu ziehen, und der Versuch, die Wurzeln in reeller Form darzustellen, f\"uhrt immer wieder auf eine cubische Gleichung von derselben Beschaffenheit. Darum hat man diesen Fall den \gesperrt{casus irreducibilis} genannt. Die cubischen Gleichungen des casus irreducibilis sind ein specieller Fall der \gesperrt{cyklischen Gleichungen mit reellen Wurzeln}, die wir im \S.\,165 kennen gelernt haben, bei deren L\"osung gleichfalls Wurzeln aus imagin\"aren Gr\"ossen, oder was auf dasselbe hinauskommt, Winkeltheilungen, vorkommen. Dass diese Wurzeln aus imagin\"aren Gr\"ossen oder Winkeltheilungen auf keine Weise zu vermeiden sind, k\"onnen wir jetzt beweisen.\footnote{Vgl. \gesperrt{H\"older}, Mathematische Annalen, Bd. 38; \gesperrt{Kneser}, ebendas., Bd. 41.}
